\documentclass[11pt]{article}
\usepackage{amsmath,amsthm,amssymb}
\usepackage[margin=1in]{geometry}
\usepackage{hyperref}
\usepackage{booktabs}
\usepackage{enumitem}

\newtheorem{theorem}{Theorem}
\newtheorem{lemma}[theorem]{Lemma}
\newtheorem{proposition}[theorem]{Proposition}
\newtheorem{corollary}[theorem]{Corollary}
\newtheorem{conjecture}[theorem]{Conjecture}
\theoremstyle{definition}
\newtheorem{definition}[theorem]{Definition}
\newtheorem{remark}[theorem]{Remark}
\newtheorem{observation}[theorem]{Observation}
\newtheorem{example}[theorem]{Example}

\newcommand{\Hq}{H_q}
\newcommand{\Sn}{S_n}
\newcommand{\Vlam}{V_\lambda}
\newcommand{\Bdec}{B_{\ell+1}^{(\lambda)}}
\newcommand{\Om}{\Omega}
\newcommand{\hatl}{\widehat\lambda}
\newcommand{\supp}{\mathrm{supp}}
\newcommand{\SYT}{\mathrm{SYT}}
\newcommand{\Rp}{{R'}}

\title{Extended sign-kill for $\hatl=(5,3)$: \\
       five-piece $E^+_{n-1}$ decomposition, three vanishing branches \\
       \large (closing the $|\hatl| = 8$ stratum)}
\author{Clio}
\date{14 May 2026 (very late prove session, eleventh pass)}

\begin{document}
\maketitle

\begin{abstract}
Let $\lambda = (5, 3, 1^q) \vdash n = q + 8$ with $q \ge 6$,
$\ell = q + 2 = n - 6$, and
$B = \Bdec\Vlam = \Rp_{\ell+1}\Rp_{\ell+2}\Rp_{\ell+3}\Rp_{\ell+4}\Rp_{\ell+5}\Rp_{\ell+6}\,\Vlam$
(six $\Rp$-blocks).  We prove the ``$\subseteq$'' direction of the
extended sign-kill conjecture for this family at the predicted threshold
$\tau = q + 3 - \hatl_1 - \hatl_2 = q - 5$:
\[
   B \;\subseteq\; \bigcap_{j=1}^{q-5} E_j^-\!\mid_{\Vlam}.
\]
At $q \in \{0, 1, 2, 3, 4, 5\}$ the predicted threshold is non-positive,
so the containment is vacuous.

The shape $\hatl = (5, 3)$ is the first case in the extended sign-kill
programme where the inner $E^+_{n-1}$ decomposition has \emph{five}
pieces (three from corner-pairs plus \emph{two} same-row pieces).  Both
same-row candidates are valid: the one at $c_1 = (1, 5)$ leaves a
non-hook inner shape $(3, 3, 1^q)$, and the one at $c_2 = (2, 3)$
leaves a hook $(5, 1^{q+1})$.  Three of the five pieces vanish in the
pullback ($\mu_{12}, \mu_h^{(1)}$ via the $\hatl{=}(4,2)$ and
$(3,3)$ recursive inputs, $\mu_h^{(2)}$ via the hook column-$1$
theorem), leaving two non-zero contributions ($\mu_{13}$ via the
$\hatl{=}(4,3)$ recursive input from the eighth pass, $\mu_{23}$ via
the $\hatl{=}(5,2)$ recursive input from the tenth pass).  The
dimensions match: $\dim B = f^{(4,2)} = 9 = f^{(3,2)} + f^{(4,1)} =
\dim B^{(\mu_{13})}_{\ell+1} + \dim B^{(\mu_{23})}_{\ell+1}$.

This completes the $|\hatl| = 8$ row of the recursive ladder (together
with the ninth-pass $(4, 4)$ paper), and is the first case where
\emph{both} recursive contributing inputs come from $|\hatl_X| = 7$
papers from this same evening.
\end{abstract}

\section{Setup}\label{sec:setup}

We use the conventions of the May-13--May-14 series.  $\Hq(\Sn)$ is
the Iwahori--Hecke algebra over $\mathbb{Q}(q)$ with generators
$T_1, \ldots, T_{n-1}$ and quadratic relation $T_j^2 = (q-1)T_j + q$.
$\Vlam$ is the Specht module in the Hoefsmit seminormal basis
$\{v_T : T \in \SYT(\lambda)\}$.

For $1 \le j \le n-1$, set $S_j := T_j + 1$.  Write
\[
   E_j^+ \;:=\; \ker(T_j - q)\!\mid_{\Vlam}
   \;=\; \mathrm{Im}(S_j\!\mid_{\Vlam}),
   \qquad
   E_j^- \;:=\; \ker(S_j)\!\mid_{\Vlam}.
\]
For $k \ge 1$, let $\Rp_k := S_{k-1}\,S_{k-2}\cdots S_1$ with the
convention $\Rp_1 = 1$.  For $\lambda \vdash n$ with
$\ell = \ell(\lambda)$, set
\[
   \Om \;=\; \Om^{(\lambda)} \;:=\; \Rp_{\ell+1}\Rp_{\ell+2}\cdots \Rp_n,
   \qquad
   B \;:=\; \Om \cdot \Vlam \;=\; \Bdec \Vlam.
\]

For $\lambda = (5, 3, 1^q)$, $n = q + 8$ and $\ell = q + 2$, so
$\Om = \Rp_{q+3}\Rp_{q+4}\Rp_{q+5}\Rp_{q+6}\Rp_{q+7}\Rp_{q+8} =
\Rp_{n-5}\Rp_{n-4}\Rp_{n-3}\Rp_{n-2}\Rp_{n-1}\Rp_n$
has \emph{six} $\Rp$-blocks.  By the May-13 multi-corner-dimension
paper, $\dim B = f^{\lambda^{(\ge 2)}} = f^{(4, 2)} = 9$ in the stable
regime $q \ge 2$.

\paragraph{Removable corners of $\lambda$.}
The removable corners are
\[
   c_1 = (1, 5), \qquad c_2 = (2, 3), \qquad c_3 = (q+2, 1).
\]
$\lambda$ has two length-preserving corners ($c_1, c_2$) and one
column-$1$ bottom corner ($c_3$), an $r = 2$ shape.

\paragraph{Identity prefix and seminormal support.}  For an SYT $T$
of any partition $\mu$, the \emph{identity prefix length} is
\[
   p(T) \;:=\; \max\{\,P \ge 0 : T(i, 1) = i \text{ for all }
   1 \le i \le P\,\}.
\]
For $W \subseteq V_\mu$,
$\supp(W) := \{\,T \in \SYT(\mu) : \exists w \in W,\ [v_T]\,w \ne 0\,\}$.

\section{Input from prior papers}\label{sec:input}

We collect the prior extended sign-kill theorems used as recursive
inputs.

\begin{proposition}[Support rule, May-14 fourth pass]\label{prop:support-rule}
Let $W \subseteq V_\nu$ be any subspace, $1 \le j \le |\nu|-1$.
If every $T \in \supp(W)$ has $j$ and $j+1$ in the same column of $T$,
then $W \subseteq E_j^-\!\mid_{V_\nu}$.
\end{proposition}

\begin{remark}\label{rem:prefix-implies-Ej-}
If every $T \in \supp(W)$ has $p(T) \ge P$, then for every
$1 \le j \le P - 1$ the entries $j, j+1$ sit at column-$1$ rows
$j, j+1$, sharing column~$1$.  By Proposition~\ref{prop:support-rule},
$W \subseteq E_j^-$ for $1 \le j \le P - 1$.
\end{remark}

\begin{theorem}[Extended sign-kill at $\hatl = (3, 3)$, May-14 seventh pass]\label{thm:33-prefix}
For $\lambda' = (3, 3, 1^{q'}) \vdash n' = q' + 6$ with $q' \ge 4$,
every $T' \in \supp\bigl(B^{(\lambda')}_{q'+3} V_{\lambda'}\bigr)$
satisfies $p(T') \ge n' - 8 = q' - 2$.  In particular
$B^{(\lambda')}_{q'+3} \subseteq E_1^-$ for $q' \ge 4$.
\end{theorem}

\begin{theorem}[Extended sign-kill at $\hatl = (4, 2)$, May-14 sixth pass]\label{thm:42-prefix}
For $\lambda' = (4, 2, 1^{q'}) \vdash n' = q' + 6$ with $q' \ge 4$,
every $T' \in \supp\bigl(B^{(\lambda')}_{q'+3} V_{\lambda'}\bigr)$
satisfies $p(T') \ge n' - 8 = q' - 2$.  In particular
$B^{(\lambda')}_{q'+3} \subseteq E_1^-$ for $q' \ge 4$.
\end{theorem}

\begin{theorem}[Extended sign-kill at $\hatl = (4, 3)$, May-14 eighth pass]\label{thm:43-prefix}
For $\lambda' = (4, 3, 1^{q'}) \vdash n' = q' + 7$ with $q' \ge 5$,
every $T' \in \supp\bigl(B^{(\lambda')}_{q'+3} V_{\lambda'}\bigr)$
satisfies $p(T') \ge n' - 10 = q' - 3$.  Consequently,
$B^{(\lambda')}_{q'+3} V_{\lambda'} \subseteq E_j^-\!\mid_{V_{\lambda'}}$
for $1 \le j \le q' - 4$.
\end{theorem}

\begin{theorem}[Extended sign-kill at $\hatl = (5, 2)$, May-14 tenth pass]\label{thm:52-prefix}
For $\lambda' = (5, 2, 1^{q'}) \vdash n' = q' + 7$ with $q' \ge 5$,
every $T' \in \supp\bigl(B^{(\lambda')}_{q'+3} V_{\lambda'}\bigr)$
satisfies $p(T') \ge n' - 10 = q' - 3$.  Consequently,
$B^{(\lambda')}_{q'+3} V_{\lambda'} \subseteq E_j^-\!\mid_{V_{\lambda'}}$
for $1 \le j \le q' - 4$.
\end{theorem}

\begin{theorem}[Hook column-$1$, May-12]\label{thm:hook-E1minus}
For $\mu = (k, 1^m) \vdash n_0$ with $k \ge 2$ and $m \ge \max(k, 2)$,
$B^{(\mu)}_{\ell_0+1} V_\mu \subseteq E_1^-\!\mid_{V_\mu}$ (where
$\ell_0 = \ell(\mu) = m + 1$).
\end{theorem}

\begin{theorem}[Phase A endpoint, May-19]\label{thm:phaseA}
For any $\Hq$-module $V$ and any $1 \le j \le n-1$,
$(T_j+1)\cdots(T_1+1) V = E_j^+(V)$.  In particular,
$\Rp_n V_\lambda = E_{n-1}^+|_{V_\lambda}$.
\end{theorem}

\begin{lemma}[Outer-factor preservation, May-14 fifth pass]\label{lem:outer-preserve}
Let $W \subseteq V_\nu$ be a subspace such that every $T \in \supp(W)$
satisfies $p(T) \ge P$.  Then for every $j \ge P + 1$,
every $T'' \in \supp((T_j{+}1) W)$ also satisfies $p(T'') \ge P$.
\end{lemma}

\section{The five-piece $E^+_{n-1}$ decomposition for $(5,3,1^q)$}\label{sec:Eplus}

The $+q$-eigenspace $E_{n-1}^+|_{V_\lambda}$ admits a clean
\emph{five}-piece decomposition reflecting the corner structure of $\lambda$.

\paragraph{Corner-pair pieces.}  For each pair $\{c_X, c_Y\}$ of distinct
removable corners, placement of $\{n-1, n\}$ at $\{c_X, c_Y\}$ in either
order yields a $T_{n-1}$-$2$-block in $V_\lambda$.  The pairs and axial
distances:

\begin{itemize}[topsep=0.2em,itemsep=0.1em]
\item $\{c_1, c_2\} = \{(1, 5), (2, 3)\}$.  Content
   $c(1, 5) - c(2, 3) = 4 - 1 = 3$.  Non-degenerate (\(|d| = 3 \ge 2\)).
   Inner shape $\mu_{12} := (4, 2, 1^q)$.
\item $\{c_1, c_3\} = \{(1, 5), (q+2, 1)\}$.  Content $c(1, 5) -
   c(q+2, 1) = 4 - (1 - (q+1)) = q + 4$.  Non-degenerate
   (\(|d| \ge 10\) at \(q \ge 6\)).  Inner shape $\mu_{13} := (4, 3, 1^{q-1})$.
\item $\{c_2, c_3\} = \{(2, 3), (q+2, 1)\}$.  Content $c(2, 3) -
   c(q+2, 1) = 1 - (1 - (q+1)) = q + 1$.  Non-degenerate
   (\(|d| \ge 7\) at \(q \ge 6\)).  Inner shape $\mu_{23} := (5, 2, 1^{q-1})$.
\end{itemize}

\paragraph{Same-row pieces.}  A same-row $+q$-piece exists for
$T_{n-1}$ on $V_\lambda$ when there is a removable corner $c$ and an
SYT of $\lambda$ with $n-1$ at the cell immediately left of $c$ (same
row) and $n$ at $c$.  The candidates:

\begin{itemize}[topsep=0.2em,itemsep=0.1em]
\item At $c_1 = (1, 5)$: requires $T(1, 4) = n-1$, $T(1, 5) = n$.
   Cells $(1, 4)$ and $(1, 5)$ sit alone in their columns
   (since $\lambda_2 = 3 < 4$, neither $(2, 4)$ nor $(2, 5)$ exists),
   so no column-strict constraint below them.  The remaining $n-2$
   entries fill $\lambda \setminus \{(1,4), (1,5)\}$, which is the
   shape $\mu_h^{(1)} := (3, 3, 1^q)$, a valid partition.
   \emph{Same-row piece at $c_1$ with inner shape $\mu_h^{(1)}$.}
\item At $c_2 = (2, 3)$: requires $T(2, 2) = n-1$, $T(2, 3) = n$.
   The remaining entries fill $\lambda \setminus \{(2,2), (2,3)\} =
   (5, 1, 1^q) = (5, 1^{q+1})$, a valid hook.  Column-$2$ constraint
   $T(1, 2) < T(2, 2) = n - 1$ requires $T(1, 2) \le n - 2$, automatic.
   Column-$3$ constraint $T(1, 3) < T(2, 3) = n$ requires
   $T(1, 3) \le n - 1$, automatic.
   \emph{Same-row piece at $c_2$ with inner hook $\mu_h^{(2)}$.}
\item At $c_3 = (q+2, 1)$: no cell to its left.
   \emph{No same-row piece at $c_3$.}
\end{itemize}

\paragraph{No other $+q$-pieces.}  Same-column pairs at column $1$ give
$-1$-eigenvectors of $T_{n-1}$ (not in $E_{n-1}^+$).

\begin{lemma}[Five-piece decomposition of $E_{n-1}^+$]\label{lem:Eplus-decomp}
Define embeddings $\Phi_X : V_{\mu_X} \hookrightarrow V_\lambda$ for
$X \in \{12, 13, 23\}$ on the seminormal basis by
\[
   \Phi_X(v_{T'}) \;:=\; v_{T_X^+} + \tau_X\,v_{T_X^-},
\]
where $T_X^+, T_X^-$ are the two SYTs of $\lambda$ obtained from
$T' \in \SYT(\mu_X)$ by placing $\{n-1, n\}$ at the pair of cells
$\{c_X, c_Y\}$ in the two possible orders, and
$\tau_X \in \mathbb{Q}(q)^\times$ is the Hoefsmit 2-block coefficient
at the corresponding axial distance.  Define
$\Phi_h^{(1)} : V_{\mu_h^{(1)}} \hookrightarrow V_\lambda$ by
$\Phi_h^{(1)}(v_{T''}) := v_{T'''}$, the unique SYT of $\lambda$ with
$n-1$ at $(1, 4)$ and $n$ at $(1, 5)$.  Define
$\Phi_h^{(2)} : V_{\mu_h^{(2)}} \hookrightarrow V_\lambda$ analogously
with $n-1$ at $(2, 2)$ and $n$ at $(2, 3)$.  Then:

\begin{enumerate}[topsep=0.3em,itemsep=0.1em,leftmargin=2em,
                  label=(\roman*)]
\item Each $\Phi_X$ ($X \in \{12, 13, 23, h^{(1)}, h^{(2)}\}$) is
   injective and $T_i$-equivariant for $1 \le i \le n - 3$.
\item The five images in $V_\lambda$ have pairwise disjoint
   seminormal supports and satisfy
   \[
       \Rp_n V_\lambda \;=\; E_{n-1}^+|_{V_\lambda}
       \;=\; \Phi_{12}(V_{\mu_{12}}) \oplus \Phi_{13}(V_{\mu_{13}})
        \oplus \Phi_{23}(V_{\mu_{23}})
        \oplus \Phi_h^{(1)}(V_{\mu_h^{(1)}})
        \oplus \Phi_h^{(2)}(V_{\mu_h^{(2)}}).
   \]
\end{enumerate}
\end{lemma}

\begin{proof}
(i) Injectivity and $T_i$-equivariance for $i \le n - 3$ are standard
(May-20 Lemma 3.2; May-14 seventh-pass Lemma~3.2): for $i \le n - 3$,
the swap $s_i$ moves only entries $i, i+1 \le n - 2$, both inside the
``inner'' cells of $\mu_X$; Hoefsmit coefficients depend only on inner
cells, so $T_i$ commutes with the corresponding $\Phi$.

(ii) Pairwise disjointness of supports: each $\Phi$'s image is
supported on SYTs of $\lambda$ where $\{n - 1, n\}$ occupies a specific
unordered pair of cells:
\begin{itemize}[topsep=0.2em,itemsep=0.1em]
\item $\Phi_{12}$: $\{(1, 5), (2, 3)\}$.
\item $\Phi_{13}$: $\{(1, 5), (q+2, 1)\}$.
\item $\Phi_{23}$: $\{(2, 3), (q+2, 1)\}$.
\item $\Phi_h^{(1)}$: $\{(1, 4), (1, 5)\}$ (with $n-1$ at $(1,4)$, $n$ at $(1, 5)$).
\item $\Phi_h^{(2)}$: $\{(2, 2), (2, 3)\}$ (with $n-1$ at $(2,2)$, $n$ at $(2, 3)$).
\end{itemize}
The five unordered cell pairs are pairwise distinct, so the supports
are pairwise disjoint.

Dimension: by Theorem~\ref{thm:phaseA}, $\Rp_n V_\lambda
= E_{n-1}^+|_{V_\lambda}$.  A direct enumeration of SYTs of $\lambda$
where $\{n-1, n\}$ lies at one of the five admissible unordered
cell-pairs above gives
$\dim E_{n-1}^+ = f^{\mu_{12}} + f^{\mu_{13}} + f^{\mu_{23}} +
f^{\mu_h^{(1)}} + f^{\mu_h^{(2)}}$,
matching the summed Phi-image dimensions, so the sum is direct.

[Computationally verified at $q \in \{2, 3, 4, 5\}$; see
\texttt{\textasciitilde/projects/scratch/2026-05-14-extended-sign-kill-53/verify\_decomp.py}.]
\qed
\end{proof}

\begin{remark}\label{rem:five-pieces}
$\hatl = (5, 3)$ is the first case in the May-14 series with a
\emph{five}-piece $E^+_{n-1}$ decomposition.  The two same-row pieces
arise because $\lambda$ has $r = 2$ (two length-preserving corners
$c_1, c_2$) \emph{and} both same-row candidates survive the
column-strict feasibility check.  For $\hatl = (4, 3)$ (eighth pass),
the same-row at $c_1 = (1, 4)$ was forbidden because cell $(2, 3)$
existed in $\lambda$, forcing $T(2, 3) > n-1$.  Here, with $\lambda_1
= 5$, the relevant blocking cell would be $(2, 4)$, which does
\emph{not} exist (since $\lambda_2 = 3$); so the constraint is vacuous
and the same-row at $c_1$ contributes.  Similarly, both same-row
pieces in $(5, 2)$ (tenth pass; only one survived there because
$\lambda_2 = 2$ made $c_2 = (2, 2)$ produce a column-$1$ conflict).
\end{remark}

\section{The two-branch reduction}\label{sec:reduction}

\begin{theorem}[Two-branch reduction]\label{thm:reduction}
For $\lambda = (5, 3, 1^q)$ with $q \ge 4$,
\begin{equation}\label{eq:reduction}
   B \;=\; (T_{n-6}{+}1)(T_{n-5}{+}1)(T_{n-4}{+}1)(T_{n-3}{+}1)(T_{n-2}{+}1)\,
   \Bigl[\,\Phi_{13}\bigl(B_{\ell(\mu_{13})+1}^{(\mu_{13})} V_{\mu_{13}}\bigr)
        + \Phi_{23}\bigl(B_{\ell(\mu_{23})+1}^{(\mu_{23})} V_{\mu_{23}}\bigr)
   \,\Bigr].
\end{equation}
\end{theorem}

\begin{proof}
We follow the May-15/May-20 pull-through, applied separately to each
of the five pieces of $\Rp_n V_\lambda$ from
Lemma~\ref{lem:Eplus-decomp}, then show three pieces vanish.

\medskip\textbf{Step 1 (pulling five $\Rp$-blocks through $\Phi$).}
Let $\Phi \in \{\Phi_{12}, \Phi_{13}, \Phi_{23}, \Phi_h^{(1)},
\Phi_h^{(2)}\}$ and $\mu$ the corresponding inner shape.

Decompose $\Rp_{n-1} = (T_{n-2}+1)\,U_1$ where $U_1 = (T_{n-3}+1)\cdots
(T_1+1) = \Rp_{n-2}^{(\mu)} \in \Hq(S_{n-2})$.  Since $T_1, \ldots,
T_{n-3}$ commute with $T_{n-1}$ (index gaps $\ge 2$), $U_1$ preserves
$E_{n-1}^+|_{V_\lambda}$.  By Lemma~\ref{lem:Eplus-decomp}(i),
$U_1 \Phi(v) = \Phi(\Rp_{n-2}^{(\mu)} v)$ for $v \in V_\mu$.  Hence
\[
   \Rp_{n-1}\Phi(V_\mu)
   \;=\; (T_{n-2}+1)\,\Phi\bigl(\Rp_{n-2}^{(\mu)} V_\mu\bigr).
\]
Iterating for $\Rp_{n-2}, \Rp_{n-3}, \Rp_{n-4}, \Rp_{n-5}$, using
that each pulled-out $T_j$ factor commutes with the next inner $\Rp$
(index gaps $\ge 2$):
\begin{equation}\label{eq:branchpull}
   \Rp_{n-5}\Rp_{n-4}\Rp_{n-3}\Rp_{n-2}\Rp_{n-1}\Phi(V_\mu)
   \;=\; \bigl(\textstyle\prod_{j=n-6}^{n-2}(T_j+1)\bigr)\,
   \Phi\bigl(Q_\mu V_\mu\bigr),
\end{equation}
where the product is in increasing-index order from left to right
($(T_{n-6}+1)$ leftmost, $(T_{n-2}+1)$ rightmost), and
\[
   Q_\mu \;:=\; \Rp_{n-6}^{(\mu)}\,\Rp_{n-5}^{(\mu)}\,\Rp_{n-4}^{(\mu)}\,
                \Rp_{n-3}^{(\mu)}\,\Rp_{n-2}^{(\mu)}.
\]

\medskip\textbf{Step 2 (combining and applying Phase A).}
By Lemma~\ref{lem:Eplus-decomp}(ii) and Theorem~\ref{thm:phaseA},
\begin{align*}
   B \;=\; \Om V_\lambda
   &\;=\; \Rp_{n-5}\Rp_{n-4}\Rp_{n-3}\Rp_{n-2}\Rp_{n-1} \Rp_n V_\lambda \\
   &\;=\; \Rp_{n-5}\Rp_{n-4}\Rp_{n-3}\Rp_{n-2}\Rp_{n-1}
      \!\!\!\bigoplus_{X}\!\!\Phi_X(V_{\mu_X}) \\
   &\;=\; \bigl(\textstyle\prod_{j=n-6}^{n-2}(T_j+1)\bigr)\,
        \sum_{X}\Phi_X\bigl(I_X\bigr),
\end{align*}
where $X$ ranges over $\{12, 13, 23, h^{(1)}, h^{(2)}\}$ and
$I_X := Q_{\mu_X} V_{\mu_X}$.

\medskip\textbf{Step 3 (the $\mu_{12}$ piece vanishes).}
$\mu_{12} = (4, 2, 1^q)$ has $\ell(\mu_{12}) = q + 2 = n - 6$, hence
$\ell(\mu_{12}) + 1 = n - 5$.  Then
$\Om^{(\mu_{12})} = \Rp_{n-5}^{(\mu_{12})}\Rp_{n-4}^{(\mu_{12})}
\Rp_{n-3}^{(\mu_{12})}\Rp_{n-2}^{(\mu_{12})}$ has four blocks, and
\[
   I_{12} \;=\; Q_{\mu_{12}} V_{\mu_{12}}
   \;=\; \Rp_{n-6}^{(\mu_{12})}\,\bigl(\Om^{(\mu_{12})} V_{\mu_{12}}\bigr)
   \;=\; \Rp_{n-6}^{(\mu_{12})}\,B^{(\mu_{12})}_{\ell(\mu_{12})+1} V_{\mu_{12}}.
\]
By Theorem~\ref{thm:42-prefix} applied to $\mu_{12} = (4, 2, 1^q)$
with $q' = q \ge 4$ (joint hypothesis of the theorem and the present
proof; we will tighten to $q \ge 6$ below),
$B^{(\mu_{12})}_{\ell+1} V_{\mu_{12}} \subseteq E_1^-|_{V_{\mu_{12}}}$.
The rightmost factor of $\Rp_{n-6}^{(\mu_{12})}$ is $(T_1+1)$, which
annihilates $E_1^-|_{V_{\mu_{12}}}$.  Hence $I_{12} = 0$.

\medskip\textbf{Step 4 (the $\mu_h^{(1)}$ piece vanishes).}
$\mu_h^{(1)} = (3, 3, 1^q)$ has $\ell(\mu_h^{(1)}) = q + 2 = n - 6$,
hence $\ell + 1 = n - 5$.  Same structure as Step~3:
\[
   I_h^{(1)} \;=\; \Rp_{n-6}^{(\mu_h^{(1)})}\,
   B^{(\mu_h^{(1)})}_{\ell+1} V_{\mu_h^{(1)}}.
\]
By Theorem~\ref{thm:33-prefix} applied to $\mu_h^{(1)} = (3, 3, 1^q)$
with $q' = q \ge 4$, $B^{(\mu_h^{(1)})}_{\ell+1} V_{\mu_h^{(1)}}
\subseteq E_1^-|_{V_{\mu_h^{(1)}}}$.  The rightmost $(T_1+1)$ of
$\Rp_{n-6}^{(\mu_h^{(1)})}$ kills it, so $I_h^{(1)} = 0$.

\medskip\textbf{Step 5 (the $\mu_h^{(2)}$ piece vanishes).}
$\mu_h^{(2)} = (5, 1^{q+1})$ has $\ell(\mu_h^{(2)}) = q + 2 = n - 6$,
hence $\ell + 1 = n - 5$.  Same structure:
\[
   I_h^{(2)} \;=\; \Rp_{n-6}^{(\mu_h^{(2)})}\,
   B^{(\mu_h^{(2)})}_{\ell+1} V_{\mu_h^{(2)}}.
\]
By Theorem~\ref{thm:hook-E1minus} applied to $\mu_h^{(2)} =
(5, 1^{q+1})$ with $k = 5$, $m = q + 1$ (hypothesis $m \ge \max(k, 2)
= 5$, i.e., $q \ge 4$ $\checkmark$), $B^{(\mu_h^{(2)})}_{\ell+1}
V_{\mu_h^{(2)}} \subseteq E_1^-|_{V_{\mu_h^{(2)}}}$.  The rightmost
$(T_1+1)$ of $\Rp_{n-6}^{(\mu_h^{(2)})}$ kills it, so $I_h^{(2)} = 0$.

\medskip\textbf{Step 6 (the $\mu_{13}$ and $\mu_{23}$ pieces are
standard $B^{(\mu_X)}$).}  Both $\mu_{13} = (4, 3, 1^{q-1})$ and
$\mu_{23} = (5, 2, 1^{q-1})$ have $\ell(\mu_X) = q + 1 = n - 7$, hence
$\ell(\mu_X) + 1 = n - 6$.  The inner quintuple product $Q_{\mu_X}$
runs from $\Rp^{(\mu_X)}_{n-6}$ to $\Rp^{(\mu_X)}_{n-2}$, exactly the
five blocks of $\Om^{(\mu_X)}$.  Therefore
\[
   I_{13} \;=\; B^{(\mu_{13})}_{\ell(\mu_{13})+1} V_{\mu_{13}},
   \qquad
   I_{23} \;=\; B^{(\mu_{23})}_{\ell(\mu_{23})+1} V_{\mu_{23}}.
\]

Combining Steps 2--6 gives~\eqref{eq:reduction}.\qed
\end{proof}

\begin{remark}\label{rem:dim-add}
The dimensions match:
\[
   \dim B \;=\; f^{(4, 2)} \;=\; 9,
   \qquad
   \dim B^{(\mu_{13})}_{\ell+1} \;=\; f^{(3, 2)} \;=\; 5,
   \qquad
   \dim B^{(\mu_{23})}_{\ell+1} \;=\; f^{(4, 1)} \;=\; 4,
\]
and $9 = 5 + 4$ confirms that both contributing Phi-pieces inject and
the five outer $(T_j+1)$ factors are jointly injective on the sum
$\Phi_{13}(I_{13}) + \Phi_{23}(I_{23})$ (computationally verified at
$q = 6$; see Section~\ref{sec:verify}).
\end{remark}

\section{Main theorem}\label{sec:main}

\begin{theorem}[Main]\label{thm:main}
For $\lambda = (5, 3, 1^q) \vdash n = q + 8$ with $q \ge 6$,
\[
   B \;\subseteq\; \bigcap_{j=1}^{q-5} E_j^-\!\mid_{V_\lambda}.
\]
\end{theorem}

\begin{proof}
By Remark~\ref{rem:prefix-implies-Ej-}, it suffices to show that every
$T \in \supp(B)$ satisfies $p(T) \ge q - 4 = n - 12$.  We use the
two-branch reduction~\eqref{eq:reduction} together with prefix bounds
on the two inner $B^{(\mu_X)}_{\ell+1}$ pieces and
Lemma~\ref{lem:outer-preserve} on the five outer factors.

\medskip\textbf{Step 1 (prefix bound on the $\Phi_{13}$-image).}
$\mu_{13} = (4, 3, 1^{q-1})$ has size $n - 2$, $\ell(\mu_{13}) =
q + 1$, and trailing-$1$-count $q'_{13} = q - 1$.  Since $q \ge 6$ we
have $q'_{13} \ge 5$, so Theorem~\ref{thm:43-prefix} applies, giving
\[
   p(T') \;\ge\; q'_{13} - 3 \;=\; q - 4
   \qquad \text{for every }
   T' \in \supp\bigl(B^{(\mu_{13})}_{\ell(\mu_{13})+1} V_{\mu_{13}}\bigr).
\]

The Phi-embedding $\Phi_{13}$ adds $\{n-1, n\}$ at $\{c_1, c_3\} =
\{(1, 5), (q+2, 1)\}$.  The added cell $(1, 5)$ sits in column $5$
(not column $1$); the added cell $(q+2, 1)$ sits at row $q + 2 = n - 6
> q + 1 = \ell(\mu_{13})$, i.e., \emph{below} all column-$1$ cells of
$\mu_{13}$.  Since $p(T') \le \ell(\mu_{13}) = q + 1$, the column-$1$
rows $1, \ldots, p(T')$ of any extension $T^{13}_\pm$ match those of
$T'$, so $p(T^{13}_\pm) \ge p(T') \ge q - 4$.

\medskip\textbf{Step 2 (prefix bound on the $\Phi_{23}$-image).}
$\mu_{23} = (5, 2, 1^{q-1})$ has size $n - 2$, $\ell(\mu_{23}) =
q + 1$, and trailing-$1$-count $q'_{23} = q - 1$.  Since $q \ge 6$ we
have $q'_{23} \ge 5$, so Theorem~\ref{thm:52-prefix} applies, giving
\[
   p(T') \;\ge\; q'_{23} - 3 \;=\; q - 4
   \qquad \text{for every }
   T' \in \supp\bigl(B^{(\mu_{23})}_{\ell(\mu_{23})+1} V_{\mu_{23}}\bigr).
\]

The Phi-embedding $\Phi_{23}$ adds $\{n-1, n\}$ at $\{c_2, c_3\} =
\{(2, 3), (q+2, 1)\}$.  Cell $(2, 3)$ is in column $3$ (not $1$); cell
$(q+2, 1)$ is in row $q + 2 > q + 1 = \ell(\mu_{23})$ (below
$\mu_{23}$'s column-$1$ cells).  By the same argument as Step~1,
$p(T^{23}_\pm) \ge p(T') \ge q - 4$.

\medskip\textbf{Step 3 (outer factors preserve the prefix).}
The five outer factors in~\eqref{eq:reduction} are $(T_j{+}1)$ for
$j \in \{n-6, n-5, n-4, n-3, n-2\} = \{q+2, q+3, q+4, q+5, q+6\}$.
For prefix bound $P = q - 4$, the hypothesis $j \ge P + 1 = q - 3$ of
Lemma~\ref{lem:outer-preserve} requires $j \ge q - 3$.  The smallest
index $j = q + 2$ satisfies $q + 2 \ge q - 3$ $\checkmark$
(unconditionally).  Applying the lemma five times in sequence, the
prefix bound $\ge q - 4$ is preserved.

\medskip\textbf{Conclusion.}  By Steps 1--3 and~\eqref{eq:reduction},
every $T \in \supp(B)$ has $p(T) \ge q - 4$.  By
Remark~\ref{rem:prefix-implies-Ej-}, $B \subseteq E_j^-$ for $1 \le
j \le q - 5 = \tau(\lambda)$.\qed
\end{proof}

\begin{remark}[Hypotheses]\label{rem:hyp}
The hypothesis $q \ge 6$ is the joint requirement of:
\begin{itemize}[topsep=0.2em,itemsep=0.1em]
\item Theorem~\ref{thm:43-prefix} at $q' = q - 1 \ge 5$, i.e., $q \ge 6$
   (for $\mu_{13}$).
\item Theorem~\ref{thm:52-prefix} at $q' = q - 1 \ge 5$, i.e., $q \ge 6$
   (for $\mu_{23}$).
\item Theorem~\ref{thm:42-prefix} at $q' = q \ge 4$ (for $\mu_{12}$
   vanishing in Step~3 of Theorem~\ref{thm:reduction}).
\item Theorem~\ref{thm:33-prefix} at $q' = q \ge 4$ (for $\mu_h^{(1)}$
   vanishing in Step~4).
\item Theorem~\ref{thm:hook-E1minus} at $\mu_h^{(2)} = (5, 1^{q+1})$
   with $k = 5$, $m = q+1 \ge 5$, i.e., $q \ge 4$ (Step~5).
\end{itemize}
The bottleneck is $q \ge 6$ from the two contributing branches.  This
is tight: at $q = 5$ the prediction $\tau = q - 5 = 0$ is vacuous, and
at $q \le 4$ the prediction is negative.  Thus $q \ge 6$ matches the
first $q$ at which $\tau \ge 1$.
\end{remark}

\section{Sharpness (computational)}\label{sec:sharp}

The predicted threshold $\tau(\hatl) = q + 3 - \hatl_1 - \hatl_2 =
q - 5$ asserts both $B \subseteq \bigcap_{j=1}^{q-5} E_j^-$
(Theorem~\ref{thm:main}) and $B \not\subseteq E_{q-4}^-$.  We verify
the latter computationally for the first non-vacuous case $q = 6$
(where $q - 4 = 2$).

\begin{theorem}[Sharpness at $q = 6$, computational]\label{thm:sharp-q6}
For $\lambda = (5, 3, 1^6)$ at $n = 14$, there exists $T \in \supp(B)$
with $p(T) = 2 = q - 4$ exactly.  Consequently $B \not\subseteq
E_{q-4}^-|_{V_\lambda} = E_2^-|_{V_\lambda}$.
\end{theorem}

\begin{proof}[Computational evidence]
Applying $\Om$ to the first $60$ Hoefsmit basis vectors at $q = 6$
gives a probe set in $\supp(B)$ of size $3771$, with prefix histogram
$\{2: 1168,\ 3: 1163,\ 4: 790,\ 5: 422,\ 6: 180,\ 7: 39,\ 8: 9\}$
(min prefix $= 2 = q - 4$, matching the predicted threshold).

An explicit witness with $p(T) = 2$ is
\[
   T(1, \cdot) = (1, 3, 5, 9, 10),
   \quad T(2, \cdot) = (2, 8, 11),
   \quad T(3..8, 1) = (4, 6, 7, 12, 13, 14).
\]
Column-$1$ entries are $(1, 2, 4, 6, 7, 12, 13, 14)$, so $p(T) = 2$
exactly (entries $1, 2$ are at the column-$1$ rows $1, 2$, but
$T(3, 1) = 4 \ne 3$, breaking the prefix).  Hence $T \in \supp(B)$
but $T$'s entries $2$ and $3$ do not share column $1$ (entry $3$ is
at $(1, 2)$), so $T \notin E_2^-|_{V_\lambda}$ (by
Proposition~\ref{prop:support-rule}'s contrapositive at $j = 2$),
hence $B \not\subseteq E_2^-$.
\qed
\end{proof}

\begin{remark}\label{rem:structural-sharp}
A structural-sharpness argument analogous to the May-14 sixth-pass
$(4, 2)$ Lemma~5.1 should locate entry $n - 10$ at $(1, 3)$ in a
forced position and produce an off-diagonal $T_{q-4}$-swap to a
tableau violating the prefix bound, contradicting $B \subseteq
E_{q-4}^-$.  The detailed write-up is omitted; the mechanism is
identical to the $(4, 2)$ and $(5, 2)$ proofs.
\end{remark}

\section{Computational verification}\label{sec:verify}

We verify Lemma~\ref{lem:Eplus-decomp} and Theorem~\ref{thm:main} at
$q = 6$ (so $n = 14$, $\dim V_\lambda = 14014$).

\paragraph{Decomposition.}  At small $q$ (where $\dim V_\lambda$ is
tractable), the 5-piece decomposition matches $\dim \Rp_n V_\lambda$:
\begin{center}
\begin{tabular}{c|c|ccccc|c|c}
\toprule
$q$ & $\dim V_\lambda$ & $f^{\mu_{12}}$ & $f^{\mu_{13}}$ &
$f^{\mu_{23}}$ & $f^{\mu_h^{(1)}}$ & $f^{\mu_h^{(2)}}$ & sum &
$\dim \Rp_n V_\lambda$ \\
\midrule
$2$ & $567$  & $90$  & $70$  & $64$  & $56$  & $35$  & $315$  & $315$  \\
$3$ & $1540$ & $189$ & $216$ & $189$ & $120$ & $70$  & $784$  & $784$  \\
$4$ & $3564$ & $350$ & $525$ & $448$ & $225$ & $126$ & $1674$ & $1674$ \\
\bottomrule
\end{tabular}
\end{center}
(Direct $\dim \Rp_n V_\lambda$ verification at larger $q$ requires
dense $N \times N$ rank computation, which is impractical at $q = 6$
where $N = 14014$; the sparse $E_1^-$ verification of the
main theorem suffices for the conclusion.)

\paragraph{$E_1^-$ containment at $q = 6$.}  Apply $\Om$ to three
independent random vectors in $V_\lambda$ (mod $P = 100\,003$ at
specialization $q = 11$, $n = 14$, $N = \dim V_\lambda = 14014$) and
check $(T_1 + 1)\,\Om v = 0$.  In each test $\Om v$ has $3771$
nonzero entries and $(T_1 + 1)\,\Om v$ is the zero vector.  Probing
$\dim B$ via $\Om$ applied to the first $60$ basis vectors gives rank
$9$, matching $f^{(4, 2)} = 9$.

\paragraph{Prefix bound and sharpness witness.}  The same probe gives
prefix histogram
$\{2: 1168,\ 3: 1163,\ 4: 790,\ 5: 422,\ 6: 180,\ 7: 39,\ 8: 9\}$,
with min $= 2 = q - 4$, matching the predicted threshold and
confirming sharpness at $q = 6$ (Theorem~\ref{thm:sharp-q6}).

Scripts:
\texttt{\textasciitilde/projects/scratch/2026-05-14-extended-sign-kill-53/verify\_decomp.py,
verify\_q6\_sparse.py}.

\section{Updated table of proved $\hatl$ rows}\label{sec:table}

\begin{table}[h]
\centering
\renewcommand{\arraystretch}{1.15}
\caption{Status of the extended sign-kill conjecture by $\hatl$.  The
$\tau$ column lists the predicted threshold $\tau(\hatl) = q + 3 -
\hatl_1 - \hatl_2$.}
\begin{tabular}{l|c|l|l}
\toprule
$\hatl$ & predicted $\tau$ & status & paper \\
\midrule
$(2, 2)$ & $q - 1$ & proved $q \ge 2$           & May-14 fourth pass \\
$(3, 2)$ & $q - 2$ & proved $q \ge 3$           & May-14 fifth pass \\
$(4, 2)$ & $q - 3$ & proved $q \ge 4$ (gap closable) & May-14 sixth pass \\
$(3, 3)$ & $q - 3$ & proved $q \ge 4$           & May-14 seventh pass \\
$(4, 3)$ & $q - 4$ & proved $q \ge 5$           & May-14 eighth pass \\
$(4, 4)$ & $q - 5$ & proved $q \ge 6$           & May-14 ninth pass \\
$(5, 2)$ & $q - 4$ & proved $q \ge 5$           & May-14 tenth pass \\
$(5, 3)$ & $q - 5$ & \textbf{proved $q \ge 6$ (this paper)} & May-14 eleventh pass \\
\bottomrule
\end{tabular}
\end{table}

The $|\hatl| \le 8$ stratum is now closed.

\section{Discussion}\label{sec:discussion}

\subsection{The $|\hatl| = 8$ row is now complete}

The ninth-pass $(4, 4)$ paper proved the first member of the
$|\hatl| = 8$ row at threshold $\tau = q - 5$.  This paper proves the
second and last member, $(5, 3)$, at the same threshold.  The
$|\hatl| \le 8$ stratum is now fully closed.

\subsection{First five-piece decomposition}

$(5, 3)$ is the first family in the May-14 series with a five-piece
$E^+_{n-1}$ decomposition (three corner-pair pieces + two same-row
pieces).  Both same-row pieces vanish in the pull-back (one via a
recursive $\hatl{=}(3,3)$ input, the other via hook column-$1$), so
the two-branch reduction~\eqref{eq:reduction} is unchanged in form
from the eighth- and tenth-pass papers --- only the bookkeeping at the
$E^+_{n-1}$ decomposition step has an extra branch.

\subsection{First case where \emph{both} contributing inputs are
recursive at $|\hatl_X| = 7$}

In $(4, 3)$ (eighth pass), the two contributing inputs were
$\hatl_X = (3, 3)$ and $(4, 2)$ (both $|\hatl| = 6$).  In $(5, 2)$
(tenth pass), the two contributing inputs were $\hatl_X = (4, 2)$
(non-hook, $|\hatl| = 6$) and a hook (handled by hook support).  In
$(5, 3)$, the two contributing inputs are $\hatl_X = (4, 3)$ and
$(5, 2)$ (both $|\hatl_X| = 7$, both proved this same evening).  This
is the recursion ``catching up'' to the immediately previous row of
the ladder, a structural milestone for the recursive programme.

\subsection{Toward a structural proof of the threshold formula}

The threshold $\tau(\hatl) = q + 3 - \hatl_1 - \hatl_2$ now holds
across all $|\hatl| \le 8$ rows by recursion.  A structural,
non-recursive proof would likely require:
\begin{itemize}[topsep=0.2em,itemsep=0.1em]
\item A uniform combinatorial description of the contributing branches
   (always the two corner pairs involving $c_3$, i.e., $\{c_i, c_3\}$
   for length-preserving $c_i$);
\item A uniform argument for the prefix bound on these contributing
   pieces, without invoking prior cases by induction.
\end{itemize}
The current case-by-case proof has the advantage of being airtight at
each step.  A uniform argument is the natural next-session target.

\subsection{Push status}

\textbf{72 unpushed commits} on \texttt{clio-vega/proofs} (after this
eleventh-pass commit; assuming the tenth-pass count of $71$ held).
PAT issue from May~6 persists --- please refresh.

\section{Status}\label{sec:status}

\paragraph{Established.}
\begin{itemize}[topsep=0.3em,itemsep=0.1em]
\item Lemma~\ref{lem:Eplus-decomp}: five-piece $E^+_{n-1}$
   decomposition for $\lambda = (5, 3, 1^q)$ at $q \ge 2$.
\item Theorem~\ref{thm:reduction}: two-branch reduction (two
   contributing, three vanishing) at $q \ge 4$.
\item Theorem~\ref{thm:main}: $B \subseteq E_j^-$ for $j \le q - 5$
   at $q \ge 6$.
\item Theorem~\ref{thm:sharp-q6}: sharpness at $j = q - 4$ for
   $q = 6$ (computational witness).
\end{itemize}

\paragraph{Open.}
\begin{itemize}[topsep=0.3em,itemsep=0.1em]
\item Structural (non-computational) sharpness across all $q \ge 6$.
\item Sharpness witnesses at $q \ge 7$ (consistent computationally
   up to memory limits; not run in this session).
\item A uniform/structural proof of the threshold formula
   $\tau(\hatl) = q + 3 - \hatl_1 - \hatl_2$ across all $\hatl$.
\item The next family $\hatl = (5, 4)$ (also $|\hatl| = 9$ along with
   $(6, 3)$ and the unique self-conjugate-like $(5, 4)$ at the
   $|\hatl| = 9$ row).
\end{itemize}

\end{document}
